\documentclass[12pt,a4j]{jarticle}
%---------------------------------------------------
\usepackage{graphicx}
\usepackage{amssymb}
\usepackage{amsmath}
\usepackage{ascmac}
\usepackage{setspace}
\usepackage{float}
\usepackage[dvips,usenames]{color}
\usepackage{colortbl}
\usepackage{algorithm}
\usepackage{algorithmic}
\usepackage{setspace} 
%---------------------------------------------------
\definecolor{bl}{rgb}{0.94,0.97,1}
\definecolor{gr}{rgb}{0.5,0.5,0.5}
\makeatletter
\def\section{\newpage\@startsection {section}{1}{\z@}{2.3ex plus -1ex minus -.2ex}{2.3 ex plus .2ex}{\Large\bf}}
\makeatother
%---------------------------------------------------
% ページの設定
\setlength{\textwidth}{160truemm}
\setlength{\textheight}{240truemm}
\setlength{\topmargin}{-14.5truemm}
\setlength{\oddsidemargin}{-0.5truemm}
\setlength{\headheight}{5truemm}
\setlength{\parindent}{1zw}
%---------------------------------------------------
%--------行間広め
\setstretch{1.5}
%--------ヘッダーとフッターの設定
\usepackage{fancyhdr}
\lhead{\leftmark}
\chead{}
\rhead{\rightmark}
\cfoot{\thepage}

\rfoot{}
\begin{document}
\vspace*{2cm}
\thispagestyle{empty}
\begin{spacing}{1}
\begin{center}
{\Large 明石工業高等専門学校専攻科 \\[1truecm]
理工学研究科修士論文} \\[3.5truecm]
\huge ビッグデータの始まりと終焉 \\
\LARGE The end of Big Data: a reasonable Internet of Things\\[4truecm]
\Large 00N7100000X 統計 太郎 \\
(経営システム工学専攻) \\[1truecm]
指導教員 指導祠堂
\end{center}

\newpage
\pagenumbering{roman}
\tableofcontents
\end{spacing}
\newpage
\abstract{高瀬舟は京都の高瀬川を上下する小舟である。徳川時代に京都の罪人が遠島を申し渡されると、本人の親類が牢屋敷へ呼び出されて、そこで暇乞いをすることを許された。それから罪人は高瀬舟に載せられて、大阪へ回されることであった。それを護送するのは、京都町奉行の配下にいる同心で、この同心は罪人の親類の中で、おも立った一人を大阪まで同船させることを許す慣例であった。これは上へ通った事ではないが、いわゆる大目に見るのであった、黙許であった。
}
\clearpage
\pagenumbering{arabic}
\pagestyle{fancy} %これをいれるとフッターとヘッダーが華やかに
\input{1_jyoron.tex}
\input{2_kanren.tex}

\pagestyle{plain}
\section*{謝辞}
ここに謝辞を書く．
\addcontentsline{toc}{section}{参考文献}
\bibliographystyle{jplain}
\begin{thebibliography}{3}
\bibitem{1}
谷啓（2011），
統計学的トロンボーン演奏法．
どこかの統計学論文誌Ａ, 0号, Vol.0, 11--92 
\bibitem{2}
つのだ☆ひろ（2009），
Rを使ったドラム演奏法．
どこかの統計学論文誌B, 0号, Vol.0, 11--92 
\bibitem{3}
Dan Aykroyd（2000），
Statistical American Joke．
{\it Journal of Blues Brothers},0号, Vol.0, 11--92 

\bibitem{4}
「ホームページの引用はあんまりしないほうがいいよ講座」{\it http://www.google.co.jp/}（最終アクセス　2013年1月2日）
\end{thebibliography}

\section*{付録}
\addcontentsline{toc}{section}{付録}

\end{document}